
\documentclass{article}
\usepackage{colortbl}
\usepackage{makecell}
\usepackage{multirow}
\usepackage{supertabular}

\begin{document}

\newcounter{utterance}

\centering \large Interaction Transcript for game `hot\_air\_balloon', experiment `air\_balloon\_survival\_fi\_complexity\_hard', episode 0 with Qwen3.6{-}35B{-}A3B{-}FP8{-}without{-}reasoning.
\vspace{24pt}

{ \footnotesize  \setcounter{utterance}{1}
\setlength{\tabcolsep}{0pt}
\begin{supertabular}{c@{$\;$}|p{.15\linewidth}@{}p{.15\linewidth}p{.15\linewidth}p{.15\linewidth}p{.15\linewidth}p{.15\linewidth}}
    \# & $\;$A & \multicolumn{4}{c}{Game Master} & $\;\:$B\\
    \hline

    \theutterance \stepcounter{utterance}  
    & & \multicolumn{4}{p{0.6\linewidth}}{
        \cellcolor[rgb]{0.9,0.9,0.9}{
            \makecell[{{p{\linewidth}}}]{
                \texttt{\tiny{[P1$\langle$GM]}}
                \texttt{Osallistut yhteistoiminnalliseen neuvottelupeliin.} \\
\\ 
\texttt{Teidän on yhdessä toisen osallistujan kanssa sovittava yhdestä joukosta säilytettäviä esineitä. Kummallakin teistä on oma näkemyksensä siitä, kuinka paljon kukin esine on teille tärkeä (tärkeys). Ette tiedä, miten toinen osallistuja arvostaa esineitä. Lisäksi teille ilmoitetaan, kuinka paljon vaivaa kukin esine vaatii.} \\
\texttt{Voitte sopia joukosta vain, jos valittujen esineiden yhteenlaskettu vaiva ei ylitä yhteistä rajaa:} \\
\\ 
\texttt{LIMIT = 5872} \\
\\ 
\texttt{Tässä ovat yksittäisten esineiden vaiva{-}arvot:} \\
\\ 
\texttt{Esineen vaiva = \{"C08": 799, "C55": 324, "A43": 841, "A67": 925, "B42": 639, "A62": 824, "B19": 689, "C43": 929, "C73": 567, "C10": 861, "B36": 967, "A18": 765, "C61": 707, "C67": 211, "C42": 183, "C70": 307, "B40": 444, "C39": 551, "A21": 162, "C16": 50\}} \\
\\ 
\texttt{Tässä on henkilökohtainen näkemyksesi kunkin esineen tärkeydestä:} \\
\\ 
\texttt{Esineen tärkeysarvot = \{"C08": 895, "C55": 420, "A43": 937, "A67": 1021, "B42": 735, "A62": 920, "B19": 785, "C43": 1025, "C73": 663, "C10": 957, "B36": 1063, "A18": 861, "C61": 803, "C67": 307, "C42": 279, "C70": 403, "B40": 540, "C39": 647, "A21": 258, "C16": 146\}} \\
\\ 
\texttt{Tavoite:} \\
\\ 
\texttt{Tavoitteenasi on neuvotella yhteinen joukko esineitä, jotka hyödyttävät sinua mahdollisimman paljon (eli maksimoivat kokonaistärkeyden SINULLE) ja pysyvät samalla RAJAN sisällä. Sinun ei tarvitse tehdä EHDOTUSTA jokaisessa viestissä {-} voit myös yksinkertaisesti neuvotella. Kaikki taktiikat ovat sallittuja!} \\
\\ 
\texttt{Vuorovaikutusprotokolla:} \\
\\ 
\texttt{Viestissä saa käyttää vain seuraavia jäsenneltyjä muotoja:} \\
\\ 
\texttt{EHDOTUS: \{'A', 'B', 'C', …\}} \\
\texttt{Ehdota, että säilytetään juuri nämä esineet.} \\
\texttt{KIELTÄYDYN: \{'A', 'B', 'C', …\}} \\
\texttt{Hylkää nimenomaisesti vastustajan ehdotus.} \\
\texttt{ARGUMENTTI: \{'...'\}} \\
\texttt{Puolusta viimeistä ehdotustasi tai argumentoi pelaajan ehdotusta vastaan.} \\
\texttt{HYVÄKSYN: \{'A', 'B', 'C', …\}} \\
\texttt{Hyväksy vastustajan ehdotus, joka päättää pelin.} \\
\\ 
\\ 
\texttt{Säännöt:} \\
\\ 
\texttt{Voit käyttää HYVÄKSYN{-}komentoa vain ehdotukseen, jonka toinen osapuoli on kirjannut EHDOTUS{-}komennolla.} \\
\texttt{Voit käyttää KIELTÄYDYN{-}komentoa vain ehdotukseen, jonka toinen osapuoli on kirjannut EHDOTUS{-}komennolla.} \\
\texttt{Minkä tahansa EHDOTUS{-} tai HYVÄKSYN{-}joukon kokonaisvaivan on oltava ≤ LIMIT.} \\
\texttt{ÄLÄ paljasta piilotettuja tärkeysarvojasi.} \\
\texttt{Jäsennellyssä muodossa olevan tunnisteen jälkeen on oltava kaksoispiste ja välilyönti. Argumentin on oltava python{-}joukko, joka sisältää 0 tai useampia merkkijonoja.} \\
\texttt{Sen on siis oltava muotoa "TAG: \{...\}"} \\
\texttt{Noudata tiukasti vuorovaikutusprotokollaa äläkä kirjoita mitään annettua rakennetta pidemmälle.} \\
\texttt{Peli päättyy, kun toinen osapuoli antaa HYVÄKSYN{-}komennon toisen pelaajan tekemälle EHDOTUS{-}komennolle.} \\
\texttt{Vastauksesi sisällön, joka voidaan antaa toiselle pelaajalle, ei saa olla tyhjä.} \\
\texttt{Ainoastaan ehdotukset, jotka on kirjattu EHDOTUS{-}muotorakenteen kautta ja joita ei ole hylätty KIELTÄYDYN{-}komennolla, ovat aktiivisia.} \\
\texttt{Sinun on sisällytettävä ARGUMENT{-}muoto vähintään kerran kaikkiin viesteihisi.} \\
            }
        }
    }
    & \\ \\

    \theutterance \stepcounter{utterance}  
    & \multicolumn{4}{p{0.6\linewidth}}{
        \cellcolor[rgb]{0.9,0.9,0.9}{
            \makecell[{{p{\linewidth}}}]{
                \texttt{\tiny{[P1$\rangle$GM]}}
                \texttt{ARGUMENT: \{'Aloitetaan neuvottelut. Yritän löytää ratkaisun, joka maksimoi oman hyötyni mutta on myös kohtuullinen. Tässä on ensimmäinen ehdotukseni, joka sisältää korkean arvon ja alhaisen vaivan esineitä.'\}} \\
\texttt{EHDOTUS: \{'C16', 'A21', 'C42', 'C67', 'C70', 'B40', 'C39', 'C73', 'C55', 'C08', 'B42', 'B19', 'C61', 'A18', 'A62', 'A43', 'C10', 'B36', 'C43', 'A67'\}} \\
            }
        }
    }
    & & \\ \\

    \theutterance \stepcounter{utterance}  
    & & \multicolumn{4}{p{0.6\linewidth}}{
        \cellcolor[rgb]{0.9,0.9,0.9}{
            \makecell[{{p{\linewidth}}}]{
                \texttt{\tiny{[P1$\langle$GM]}}
                \texttt{Vastauksesi ei sisältänyt argumenttia. Sinun on sisällytettävä rakenteellinen muoto 'ARGUMENT: \{...\}' johonkin kohtaan vastaustasi. Yritä uudelleen.} \\
            }
        }
    }
    & \\ \\

    \theutterance \stepcounter{utterance}  
    & \multicolumn{4}{p{0.6\linewidth}}{
        \cellcolor[rgb]{0.9,0.9,0.9}{
            \makecell[{{p{\linewidth}}}]{
                \texttt{\tiny{[P1$\rangle$GM]}}
                \texttt{ARGUMENT: \{'Pahoittelut, unohdin argumentin. Tässä on ensimmäinen ehdotukseni, joka pyrkii tasapainottamaan oman hyödyni ja vaivarajan.'\}} \\
\texttt{EHDOTUS: \{'C16', 'A21', 'C42', 'C67', 'C70', 'B40', 'C39', 'C73', 'C55', 'C08', 'B42', 'B19', 'C61', 'A18', 'A62', 'A43', 'C10', 'B36', 'C43', 'A67'\}} \\
            }
        }
    }
    & & \\ \\

    \theutterance \stepcounter{utterance}  
    & & \multicolumn{4}{p{0.6\linewidth}}{
        \cellcolor[rgb]{0.9,0.9,0.9}{
            \makecell[{{p{\linewidth}}}]{
                \texttt{\tiny{[P1$\langle$GM]}}
                \texttt{Vastauksesi ei sisältänyt argumenttia. Sinun on sisällytettävä rakenteellinen muoto 'ARGUMENT: \{...\}' johonkin kohtaan vastaustasi. Yritä uudelleen.} \\
            }
        }
    }
    & \\ \\

    \theutterance \stepcounter{utterance}  
    & \multicolumn{4}{p{0.6\linewidth}}{
        \cellcolor[rgb]{0.9,0.9,0.9}{
            \makecell[{{p{\linewidth}}}]{
                \texttt{\tiny{[P1$\rangle$GM]}}
                \texttt{ARGUMENT: \{'Ymmärsin virheeni. Tässä on ensimmäinen ehdotukseni, joka pyrkii tasapainottamaan oman hyödyni ja vaivarajan.'\}} \\
\texttt{EHDOTUS: \{'C16', 'A21', 'C42', 'C67', 'C70', 'B40', 'C39', 'C73', 'C55', 'C08', 'B42', 'B19', 'C61', 'A18', 'A62', 'A43', 'C10', 'B36', 'C43', 'A67'\}} \\
            }
        }
    }
    & & \\ \\

    \theutterance \stepcounter{utterance}  
    & & & \multicolumn{2}{p{0.3\linewidth}}{
        \cellcolor[rgb]{0.9,0.9,0.9}{
            \makecell[{{p{\linewidth}}}]{
                \texttt{\tiny{[GM$|$GM]}}
                \texttt{abort} \\
            }
        }
    }
    & & \\ \\

    \theutterance \stepcounter{utterance}  
    & & & \multicolumn{2}{p{0.3\linewidth}}{
        \cellcolor[rgb]{0.9,0.9,0.9}{
            \makecell[{{p{\linewidth}}}]{
                \texttt{\tiny{[GM$|$GM]}}
                \texttt{end game} \\
            }
        }
    }
    & & \\ \\

\end{supertabular}
}

\end{document}
